%  LaTeX support: latex@mdpi.com
%  manuscriptR1V1.tex: Round-1 revision rewrite (skeleton only).
%  Based on the MDPI Mathematics template from first_round/.

%=================================================================
\documentclass[mathematics,article,submit,moreauthors]{Definitions/mdpi}

%=================================================================
% MDPI internal commands - do not modify
\firstpage{1}
\makeatletter
\setcounter{page}{\@firstpage}
\makeatother
\pubvolume{1}
\issuenum{1}
\articlenumber{0}
\pubyear{2026}
\copyrightyear{2026}
%\externaleditor{Firstname Lastname}
\datereceived{ }
\daterevised{ }
\dateaccepted{ }
\datepublished{ }

%=================================================================
\usepackage{xcolor}

%=================================================================
% Full title of the paper (Capitalized)

\Title{Aggregation Semantics for Prioritization in Cooperative Automated Decision Making}

% Author Orchid ID: enter ID or remove command
\newcommand{\orcidauthorA}{0000-0003-3267-6753}
\newcommand{\orcidauthorB}{0009-0002-0183-7905}
\newcommand{\orcidauthorC}{0000-0002-7653-1452}

% Authors, for the paper (add full first names)
\Author{Pavel Novoa-Hern\'andez$^{1,*}$\orcidA{}, Mariia Godz$^{2}$\orcidA{} and David A. Pelta$^{2}$\orcidC{}}

% MDPI internal command: Authors, for metadata in PDF
\AuthorNames{Pavel Novoa-Hern\'andez, Mariia Godz and David A. Pelta}

% Affiliations / Addresses
\address{%
$^{1}$ \quad Dept. Computer Engineering and Systems, Universidad de La Laguna; pnovoahe@ull.edu.es\\
$^{2}$ \quad Dept. Computer Science and Artificial Intelligence, Universidad de Granada; \{mariiagodz,dpelta\}@ugr.es}

% Contact information of the corresponding author
\corres{Correspondence: pnovoahe@ull.edu.es}

% Abstract placeholder (to be rewritten with the new structure)
\abstract{To be rewritten.}

% Keywords
\keyword{
cooperative automated decision making;
aggregation operators;
multi-criteria decision making;
context-aware decision support;
prioritization
}

%%%%%%%%%%%%%%%%%%%%%%%%%%%%%%%%%%%%%%%%%%
\begin{document}

% ==============================================================================
% NOTA GLOBAL DE ESTILO:
% [Aborda 2.8 y M1.7]: Todo el documento debe ser reescrito en tercera persona
% (evitar "We provide", "We compare") y con una revisi\'on exhaustiva de ingl\'es.
% ==============================================================================

\section{Introduction}
% [Aborda 1.1 y 2.2]: Explicar claramente la novedad. No estamos inventando
% nuevos operadores, sino analizando su sem\'antica en el contexto de CADEMAS.
% [Aborda 2.3]: Formulaci\'on expl\'icita de las Research Questions (RQs).

\section{Background and Related Work}
% [Aborda 1.1 y 2.2]: Distinguir nuestra contribuci\'on del MCDM tradicional y la l\'ogica difusa.
% [Aborda 2.3]: A\~nadir art\'iculos recientes sobre sistemas de decisi\'on sensibles al contexto.
% [Aborda 2.4]: Incluir una tabla cronol\'ogica al final de la secci\'on resumiendo
% la literatura y marcando el hueco que llena nuestro paper.


This section provides the conceptual foundations required for the remainder
of the paper. We first review the literature on context-aware decision
support systems and aggregation operators, highlighting the gap between
the representation of contextual knowledge and the analysis of the
aggregation mechanisms used to combine contextual and predictive
information. We then introduce the Cooperative Automated Decision-Making
System (CADEMAS) framework and its machine-learning instantiation
(CADEMAS-ML), which provide the architectural setting in which the
aggregation strategies studied in this work are analyzed.

\subsection{Related Work}
\label{subsec:related-work}

Research on cooperative and context-aware decision support has established
multiple mechanisms for incorporating human knowledge into data-driven
systems. These include symbolic rules, ontologies, clinical pathways,
and human-centred interaction principles. For example, Teng and
Tan~\cite{TengTan2008} proposed cognitive agents that combine
propositional rules with reinforcement learning to support context-aware
decision making. In the clinical domain, Yao and
Kumar~\cite{YaoKumar2012} employed ontologies to identify
context-appropriate rules and instantiate clinical-pathway fragments for
individual patients. More recently, Cronholm and
G\"obel~\cite{CronholmGoebel2022} argued that effective decision support
requires combining human and artificial intelligence through appropriate
interaction and governance mechanisms rather than relying solely on
predictive performance. Complementing these studies, the survey of
Merkert et~al.~\cite{MerkertEtAl2015} found that the usefulness of
machine learning in decision support depends strongly on the decision
task, process stage, and organizational context, while human and
organizational evaluation criteria remain comparatively underexplored.

Collectively, these works demonstrate the importance of contextual
knowledge in decision support systems. However, their primary focus is
typically on knowledge representation, architecture design, rule
selection, or human--AI interaction. Comparatively less attention has
been paid to the numerical integration step that combines predictive
evidence with contextual assessments once both sources of information are
available. As a result, simple weighted averages are often adopted as a
convenient integration mechanism, even though their compensatory
semantics may allow strong predictive evidence to offset weak contextual
support. While contextual rules can impose constraints upstream, this
does not necessarily guarantee that such constraints are preserved after
score fusion.

The aggregation literature provides the mathematical foundations for
analyzing this issue. Among the most influential contributions, Yager's
ordered weighted averaging (OWA) operator~\cite{Yager1988} introduced a
family of aggregation functions capable of representing different
degrees of conjunctive and disjunctive behaviour through ordered
weights. Yager~\cite{Yager2005} subsequently extended this framework by
combining OWA operators with t-norms and relating them to Choquet and
Sugeno integrals. More broadly, aggregation theory has shown that the
choice of operator determines the degree of compensation allowed among
criteria and therefore influences the resulting decision semantics. This
importance is reflected in the review by Mardani
et~al.~\cite{Mardani2018}, which identified OWA-based methods as one of
the most widely used aggregation families in fuzzy MCDM applications.
Similarly, Csisz\'ar~\cite{Csiszar2021} characterized OWA as a flexible
parameterized framework spanning numerous well-known operators while
highlighting the challenge of selecting appropriate aggregation
behaviour.

A related line of research has investigated aggregation in collective
decision-making settings. Moral et~al.~\cite{MoralEtAl2017} showed that
both the aggregation operator and the associated distance measure affect
the speed and quality of consensus formation. Trillo
et~al.~\cite{TrilloEtAl2023} further explored customizable aggregation
mechanisms that allow experts to express different attitudes toward
criterion importance and consensus requirements. These studies provide
valuable insights into compensation, consensus, and preference
aggregation. Nevertheless, their focus remains on combining preferences
or judgments from multiple decision makers rather than integrating
predictive outputs from machine-learning models with contextual signals
that encode organizational priorities, regulations, or policy
constraints.

Consequently, a gap remains between context-aware decision-support
frameworks and aggregation theory. The former provides mechanisms for
representing and managing contextual knowledge but rarely examines the
behavioral implications of alternative aggregation semantics. The latter
offers a rich collection of aggregation operators and theoretical
properties but has seldom been studied from the perspective of
policy-compliant score fusion in cooperative automated decision-making
architectures. In particular, relatively little attention has been paid
to questions such as whether contextual assessments can effectively act
as veto mechanisms, how aggregation semantics influence robustness to
predictive overconfidence, or how prioritization rankings respond to
uncertainty in predictive and contextual information. These questions
motivate the analysis developed in this work. Table~\ref{tab:related-works-comparison}
summarizes representative contributions along these dimensions and
highlights the gap addressed here.

\begin{table}[t]
\caption{Chronological comparison of representative related work. ``Partially'' indicates that a contribution can encode rules, conjunctive criteria, or minimum thresholds, but does not establish veto preservation for predictive--contextual score fusion.}
\label{tab:related-works-comparison}
\renewcommand{\arraystretch}{1.25}\footnotesize
\begin{adjustwidth}{-\extralength}{0cm}
\begin{tabularx}{\fulllength}{>{\raggedright\arraybackslash}p{2.6cm}|>{\raggedright\arraybackslash}X|>{\raggedright\arraybackslash}X|C|C}
\toprule
\textbf{Reference} & \textbf{Application / Framework} & \textbf{Aggregation method} & \textbf{Strict contextual veto} & \textbf{Operator robustness} \\
\midrule
Yager~\cite{Yager2005} (2005) &
Multicriteria decision making &
TOWA (t-norm--OWA); Choquet and Sugeno extensions &
Partially & Partially \\
\midrule
Teng and Tan~\cite{TengTan2008} (2008) &
Context-aware DSS via cognitive agents &
Symbolic rules with reinforcement learning; scalar fusion not specified &
Partially & No \\
\midrule
Yao and Kumar~\cite{YaoKumar2012} (2012) &
Ontology-based clinical DSS &
Context-driven rule / pathway selection; scalar fusion not specified &
Partially & No \\
\midrule
Merkert et~al.~\cite{MerkertEtAl2015} (2015) &
Survey of ML in DSS &
Not a specific aggregation operator &
No & No \\
\midrule
Moral et~al.~\cite{MoralEtAl2017} (2017) &
Group decision making and consensus &
Aggregation operators combined with distance functions &
No & Partially \\
\midrule
Mardani et~al.~\cite{Mardani2018} (2018) &
Review of fuzzy MCDM &
Reviews OWA and related fuzzy aggregation families &
No & No \\
\midrule
Csisz\'ar~\cite{Csiszar2021} (2021) &
OWA theory and applications &
Parameterized OWA family; weight selection &
Partially & Partially \\
\midrule
Cronholm and G\"obel~\cite{CronholmGoebel2022} (2022) &
Human-centred AI decision support &
Human--AI design principles; scalar fusion not specified &
No & No \\
\midrule
Trillo et~al.~\cite{TrilloEtAl2023} (2023) &
Consensus-based multicriteria GDM &
Expert-customized operator; per-criterion minimum consensus &
Partially & No \\
\midrule
\textbf{Present work} &
\textbf{Cooperative ADM (CADEMAS)} &
\textbf{Comparative linear, geometric, and minimum fusion of predictive and contextual scores} &
\textbf{Yes} & \textbf{Yes} \\
\bottomrule
\end{tabularx}
\end{adjustwidth}
\end{table}

% [Aborda 1.1 y 2.2]: Distinguir nuestra contribuci\'on del MCDM tradicional y la l\'ogica difusa.
Unlike traditional fuzzy MCDM studies, the objective of this work is not
to propose a new aggregation operator, derive additional algebraic
properties, or solve a specific multicriteria decision problem through a
particular aggregation scheme. Instead, aggregation operators are treated
as alternative policy semantics within a cooperative automated
decision-making architecture. Likewise, unlike many context-aware DSS
frameworks, the focus is not on representing contextual knowledge itself
but on understanding how different integration mechanisms affect the
interaction between predictive evidence and contextual constraints.
Consequently, the contribution of this work lies at the architectural and
methodological level, namely, the systematic analysis of aggregation as a
governance mechanism within CADEMAS-ML.

\subsection{CADEMAS and CADEMAS-ML}
\label{sect:cademas}

The literature reviewed in the previous section shows that many
context-aware decision-support systems focus on representing contextual
knowledge through rules, ontologies, policies, or human expertise.
However, comparatively less attention has been paid to the integration
mechanism that combines contextual assessments with predictive evidence.
The CADEMAS framework introduced by \citet{NovoaHernandez2026} provides
a suitable setting for studying this issue because it explicitly
separates evidence generation from contextual evaluation and decision
integration.

CADEMAS is a general framework for designing cooperative automated
decision-making systems in which multiple heterogeneous Automated
Decision-Making (ADM) units collaborate to produce a unified
recommendation, score, or label. The framework is intentionally agnostic
with respect to the internal implementation of these units. An ADM may
correspond to a machine-learning model, a rule-based system, an
optimization algorithm, an expert system, or any other computational
mechanism capable of generating decision evidence. Likewise, contextual
knowledge may be represented through expert rules, organizational
policies, regulatory requirements, or other knowledge-representation
formalisms.

Rather than prescribing how individual decisions are produced, CADEMAS
defines an architectural process through which heterogeneous decision
signals are coordinated and subsequently evaluated against contextual
considerations before a final recommendation is generated.

Formally, a CADEMAS instance is represented by
\begin{equation}
\mathcal{S}=
\langle I, U, O, K, T \rangle,
\end{equation}
where $I$ denotes the Input Manager, $U$ the set of ADM units, $O$ the
Decision Integration Layer, $K$ the Contextual Modulator, and $T$ the
universe of admissible data representations.

The resulting architecture follows a layered decision process. First,
the ADM units independently generate decision evidence from the
available information. Different units may operate on different subsets
of the input data and employ different reasoning mechanisms. In
parallel, the Contextual Modulator evaluates the same decision case from
the perspective of organizational priorities, domain knowledge,
regulatory requirements, or other contextual considerations, producing a
degree of contextual satisfaction. Finally, the Decision Integration
Layer combines both sources of information into a unified decision
output.

This explicit separation between predictive assessment, contextual
evaluation, and decision integration is one of the defining
characteristics of CADEMAS. It allows the same architecture to be
instantiated across different domains while making the integration
mechanism itself an independent design component.

A particularly relevant specialization is CADEMAS-ML, in which the ADM
units correspond to machine-learning models. In this case, the
predictive subsystem resembles a cooperative ensemble in which multiple
models contribute to a unified predictive assessment. Unlike
conventional ensemble methods, however, CADEMAS-ML separates predictive
cooperation from contextual evaluation. Predictive outputs are first
consolidated into a common predictive score and only afterwards combined
with contextual knowledge that encodes organizational or institutional
priorities. Following previous CADEMAS-ML implementations
\citep{PerezCanedo2022,PerezCanedo2024,Godz2026}, contextual knowledge
is represented through an interpretable fuzzy inference system.

To focus on the role of aggregation independently of any specific
application domain, we abstract the predictive subsystem by a normalized
predictive score
\begin{equation}
R_i \in [0,1],
\end{equation}
which summarizes the evidence generated by the cooperating ADM units for
decision case $i$. Depending on the application, $R_i$ may represent a
risk estimate, a confidence score, a probability of occurrence, or any
other normalized predictive assessment.

Independently, the contextual subsystem produces a normalized context
score
\begin{equation}
Q_i \in [0,1],
\end{equation}
which measures the extent to which decision case $i$ satisfies the
relevant contextual requirements, policies, or priorities.

The predictive and contextual assessments are then combined through an
aggregation operator $A$
\begin{equation}
P_i = A(R_i,Q_i),
\end{equation}
yielding the final prioritization score. The resulting scores are used
to rank candidate cases and determine the corresponding Top-$K$
alternatives.

Existing CADEMAS-ML applications employ a linear convex aggregation rule
to combine predictive and contextual information. However, the CADEMAS
framework itself does not prescribe any specific aggregation semantics
\citep{NovoaHernandez2026}. Consequently, the aggregation operator can
be regarded as an independent policy-design component that determines
how predictive evidence and contextual considerations interact. This
property makes CADEMAS-ML an appropriate testbed for investigating how
alternative aggregation semantics influence prioritization behavior,
policy compliance, robustness, and ranking stability while preserving
the same predictive and contextual subsystems.


\section{Theoretical Framework and Aggregation Properties}
% Esta secci\'on sube el rigor matem\'atico y generaliza los resultados.

\subsection{Preliminaries and Assumptions}
% [Aborda 2.5 y 1.5]: Declarar la universalidad. Definir el dominio [0,1],
% y explicar qu\'e ocurre si no est\'an normalizados. Quitar "normalizaci\'on" como
% propiedad y formalizar conmutatividad e idempotencia.
% [Aborda M1.1 y M1.2]: Aclarar terminolog\'ia (t-norma vs operador agregador) y unificar notaci\'on.
% [Aborda 1.2]: Antiguo Teorema 2 se convierte en el Lema 1 (Monotonicidad y Fronteras).

% [Aborda 2.5]: Universalidad y asunciones del dominio.
Let $R \in [0,1]$ represent the predictive score and $Q \in [0,1]$ represent the contextual score assigned to a given decision alternative. While our analysis assumes these scores are normalized to the unit interval, this assumption does not restrict the universality of our findings. Any bounded continuous scoring domain $[a,b]$ can be isomorphically mapped to $[0,1]$ through standard min-max scaling. Consequently, the theoretical properties and operator behaviors discussed herein hold universally for any bounded decision space.

% [Aborda 1.5 y M1.1]: Corrección de la "normalización" y aclaración de terminología.
In this framework, the integration of predictive and contextual dimensions is governed by an \textit{aggregation function}. Formally, a bivariate aggregation function is a mapping $A: [0,1]^2 \to [0,1]$ that satisfies the following fundamental axioms:
\begin{itemize}
    \item \textit{Boundary conditions:} $A(0,0) = 0$ and $A(1,1) = 1$.
    \item \textit{Monotonicity:} $A(x_1, y_1) \le A(x_2, y_2)$ whenever $x_1 \le x_2$ and $y_1 \le y_2$.
\end{itemize}

Depending on the chosen mechanism, additional algebraic properties may be desirable, such as commutativity ($A(x,y) = A(y,x)$) and idempotency ($A(x,x) = x$). It is crucial to distinguish between general aggregation functions and specific mathematical sub-families. For example, while triangular norms (t-norms) such as the standard minimum operator are a specific class of aggregation functions that strictly require commutativity and associativity, not all operators analyzed in cooperative frameworks belong to this class.

% [Aborda 1.2]: El antiguo Teorema 2 rebajado a Lema fundacional.
Building upon these standard axioms, we establish the following baseline lemma, which governs the structural behavior of any valid aggregation process within our system.

\begin{Lemma}
    \label{lemma:monotonicity-and-boundary-preservation}
    \textbf{Lemma 1 (Monotonicity and Boundary Preservation).} \textit{For any candidate evaluated under an aggregation function $A$ satisfying the standard axioms, an isolated improvement in either the predictive score $R$ or the contextual score $Q$ guarantees that the final integrated score will not decrease. Furthermore, the evaluation is strictly bounded such that the presence of a perfect partial score cannot force the aggregation beyond the strict limits of the codomain, meaning $A(R,0) \ge 0$ and $A(R,1) \le 1$ for any $R \in [0,1]$.}
\end{Lemma}



\subsection{Choice of Operators and Archetypes}
% [Aborda 1.12 y 2.6]: Justificar por qu\'e elegimos Lineal, Geom\'etrico y M\'inimo.
% Explicar que representan los arquetipos de compensaci\'on y por qu\'e descartamos Choquet/OWA emp\'iricamente.
% [Aborda 1.4]: Especificar que para el geom\'etrico lambda $\in (0,1)$ para evitar $0^0$.

\subsection{Veto Preservation and Compensation Limits}
% [Aborda 1.2 y 1.3]: Reformular la Proposici\'on 1 para clases m\'as amplias de
% operadores usando propiedades exactas ($A(x,0)=0$ y $A(x,y)>0$ si $x,y>0$).
% [Aborda 2.1]: Establecer las condiciones matem\'aticas exactas (Corolario/Proposici\'on)
% para que ocurra el ``rank reversal'' bajo agregaci\'on lineal.

\subsection{Sensitivity and Robustness to Contextual Imprecision}
% [Aborda 1.9]: NUEVO TEOREMA 4. Formular el resultado anal\'itico evaluando
% $|A(R, Q+\varepsilon) - A(R, Q)|$ para justificar la superioridad del operador m\'inimo.

\section{Experimental Evaluation: Monte Carlo Simulations}
% Separamos las simulaciones del caso de uso real para mayor claridad.

\subsection{Experimental Protocol and Metrics}
% [Aborda 1.7]: Aumentar $N$ (ej. $N=1000$ en lugar de 30), explicar intervalos
% de confianza y reportar semillas aleatorias para reproducibilidad.
% [Aborda 1.11]: Redefinir la m\'etrica de Policy Violation $V_K$ para evitar
% divisi\'on por cero (e.g., definida a trozos).
% [Aborda M1.4 y M1.5]: Explicar c\'omo se rompen los empates (tie-breaking) en el
% operador m\'inimo y aclarar el funcionamiento de argsort (ascendente/descendente).

\subsection{Compensation Effects under Intermediate Contexts}
% [Aborda 1.6]: Mostrar experimentos donde el contexto NO es cero ($Q > 0$),
% para ver c\'omo se comportan los operadores en grises.
% [Aborda M1.3]: Definir formalmente y sin ambig\"uedad qu\'e significa
% ``predictive overconfidence'' en la simulaci\'on.

\subsection{Robustness to Contextual Uncertainty}
% [Aborda 1.8 y 2.7]: An\'alisis de Sensibilidad. Evaluar m\'ultiples valores de $\lambda$,
% diferentes modelos de ruido (no solo uno) y distribuciones de poblaci\'on.
% *Nota: Aqu\'i referenciamos el nuevo resultado anal\'itico de la Secci\'on 3.4.

\section{Real-World Case Study: Employee Attrition Prioritization}
% [Aborda 1.10]: Aclarar fuertemente que esto es una ILUSTRACI\'ON basada en
% scores precomputados de CADEMAS-ML, no una validaci\'on independiente de un modelo predictivo.

\subsection{Dataset Setup and Contextual Vetoes}
% [Aborda 1.10]: Explicar c\'omo se seleccionaron los empleados y c\'omo
% se obtuvieron los context scores.

\subsection{Prioritization Results and Top-K Stability}
% [Aborda 1.10 y 2.7]: Presentar los resultados y mostrar una gr\'afica/tabla
% de sensibilidad demostrando si las conclusiones se mantienen estables
% para distintos valores de $K$ (ej. Top-5, Top-10, Top-50).

\section{Discussion and Conclusions}
% Resumen de los hallazgos.
% Implicaciones pr\'acticas (la elecci\'on del operador es gobernanza).
% Trabajo futuro.

%%%%%%%%%%%%%%%%%%%%%%%%%%%%%%%%%%%%%%%%%%
\authorcontributions{To be completed.}

\funding{This research was supported by the project ``Study, Analysis and Evaluation of Cooperative Automated Decision-Making Systems (CADEMAS)'', PID2023-146575NB-I00, funded by MCIU/AEI/I0.13039/501100011033 and by FSE+.}

\institutionalreview{Not applicable.}

\informedconsent{Not applicable.}

\dataavailability{To be completed.}

\acknowledgments{Mariia Godz has a pre-doctoral research contract associated with the project ``Study, Analysis and Evaluation of Cooperative Automated Decision-Making Systems (CADEMAS)'', PID2023-146575NB-I00, funded by MCIU/AEI/I0.13039/501100011033 and by FSE+.}

\conflictsofinterest{The authors declare no conflicts of interest.}

%%%%%%%%%%%%%%%%%%%%%%%%%%%%%%%%%%%%%%%%%%
\begin{adjustwidth}{-\extralength}{0cm}
\reftitle{References}

\begin{thebibliography}{999}
    
    \bibitem[Kierner et~al.(2023)Kierner, Kucharski, and Kierner]{Kierner2023}
    Kierner, S.; Kucharski, J.; Kierner, Z.
    \newblock Taxonomy of Hybrid Architectures Involving Rule-Based Reasoning and
    Machine Learning in Clinical Decision Systems: A Scoping Review.
    \newblock {\em Journal of Biomedical Informatics} {\bf 2023}, {\em
        144},~104428.
    \newblock {\url{https://doi.org/10.1016/j.jbi.2023.104428}}.
    
    \bibitem[Rahman et~al.(2025)Rahman, Zainal, Mubarak, Shabir, Anwar, and
    Asrowardi]{Rahman2025}
    Rahman, T.; Zainal, D.; Mubarak, H.; Shabir, F.; Anwar, N.; Asrowardi, I.
    \newblock Utilizing Machine Learning-Based Decision-Making to Align Higher
    Education Curriculum with Industry Requirements.
    \newblock {\em International Journal of Modern Education and Computer Science}
    {\bf 2025}, {\em 17},~1--15.
    \newblock {\url{https://doi.org/10.5815/ijmecs.2025.04.01}}.
    
    \bibitem[Pérez-Domínguez et~al.(2024)Pérez-Domínguez, Callejas-Cuervo, and
    García-Luna]{PerezDominguez2024}
    Pérez-Domínguez, L.; Callejas-Cuervo, M.; García-Luna, F.
    \newblock Intuitionistic Fuzzy Dimensional Analysis in Multi-Criteria
    Decision-Making: A Computational Approach.
    \newblock {\em IEEE Access} {\bf 2024}, {\em 12},~1--15.
    \newblock {\url{https://doi.org/10.1109/ACCESS.2024.3422616}}.
    
    \bibitem[Novoa-Hernández et~al.(2026)Novoa-Hernández, Pelta, Godz, Verdegay,
    and Buendia-Carrillo]{NovoaHernandez2026}
    Novoa-Hernández, P.; Pelta, D.; Godz, M.; Verdegay, J.; Buendia-Carrillo, D.
    \newblock {CADEMAS} – A Framework for Cooperative Automated Decision-Making
    Systems.
    \newblock In Proceedings of the Proceedings of the 2026 IEEE Conference on
    Artificial Intelligence (CAI),  2026, pp. 756--761.
    \newblock {\url{https://doi.org/10.1109/CAI68641.2026.11536392}}.
    
    \bibitem[{ELI Innovation Paper}(2022)]{ELI_ADM_Principles_2022}
    {ELI Innovation Paper}.
    \newblock Guiding Principles for Automated Decision-Making in the {EU}.
    \newblock Innovation paper, European Law Institute (ELI),  2022.
    
    \bibitem[Godz et~al.(2026)Godz, Novoa-Hernández, and Pelta]{Godz2026}
    Godz, M.; Novoa-Hernández, P.; Pelta, D.
    \newblock A Cooperative and Context-Aware Approach for Employee Attrition
    Prevention. In {\em Information Processing and Management of Uncertainty in
        Knowledge-Based Systems}; Springer Nature Switzerland,  2026; pp. 203--217.
    \newblock {\url{https://doi.org/10.1007/978-3-032-29000-7_15}}.
    
    \bibitem[Yager(1988)]{Yager1988}
    Yager, R.
    \newblock On Ordered Weighted Averaging Aggregation Operators in Multicriteria
    Decisionmaking.
    \newblock {\em IEEE Transactions on Systems, Man, and Cybernetics} {\bf 1988},
    {\em 18},~183--190.
    \newblock {\url{https://doi.org/10.1109/21.87068}}.
    
    \bibitem[Tan and Chen(2010)]{Tan2009}
    Tan, C.; Chen, X.
    \newblock Induced Intuitionistic Fuzzy Ordered Weighted Geometric Operator and
    Its Application to Multiple Attribute Decision Making.
    \newblock {\em International Journal of Intelligent Systems} {\bf 2010}, {\em
        25},~1193--1209.
    \newblock {\url{https://doi.org/10.1002/int.20446}}.
    
    \bibitem[Grabisch et~al.(2000)Grabisch, Murofushi, and Sugeno]{Grabisch1996}
    Grabisch, M.; Murofushi, T.; Sugeno, M.
    \newblock Aggregation Operators and Fuzzy Integrals. In {\em Fuzzy Measures and
        Integrals: Theory and Applications}; Physica-Verlag,  2000; pp. 3--41.
    
    \bibitem[Riaz et~al.(2022)Riaz, Farid, Wang, and Pamucar]{Riaz2022}
    Riaz, M.; Farid, H.; Wang, W.; Pamucar, D.
    \newblock Interval-Valued Linear Diophantine Fuzzy Frank Aggregation Operators
    with Multi-Criteria Decision-Making.
    \newblock {\em Mathematics} {\bf 2022}, {\em 10},~1811.
    \newblock {\url{https://doi.org/10.3390/math10111811}}.
    
    \bibitem[Senapati et~al.(2022)Senapati, Chen, and Yager]{Senapati2022}
    Senapati, T.; Chen, G.; Yager, R.
    \newblock Aczel--Alsina Aggregation Operators and Their Application to
    Intuitionistic Fuzzy Multiple Attribute Decision Making.
    \newblock {\em International Journal of Intelligent Systems} {\bf 2022}, {\em
        37},~1529--1551.
    \newblock {\url{https://doi.org/10.1002/int.22684}}.
    
    \bibitem[Akram et~al.(2020)Akram, Yaqoob, Ali, and Chammam]{Akram2020}
    Akram, M.; Yaqoob, N.; Ali, G.; Chammam, W.
    \newblock Extensions of Dombi Aggregation Operators for Decision Making under
    m-Polar Fuzzy Information.
    \newblock {\em Journal of Mathematics} {\bf 2020}, {\em 2020},~4739567.
    \newblock {\url{https://doi.org/10.1155/2020/4739567}}.
    
    \bibitem[Mardani et~al.(2018)Mardani, Nilashi, Zavadskas, Awang, Zare, and
    Jamal]{Mardani2018}
    Mardani, A.; Nilashi, M.; Zavadskas, E.K.; Awang, S.; Zare, H.; Jamal, N.M.
    \newblock Decision making methods based on fuzzy aggregation operators: Three
    decades review from 1986 to 2017.
    \newblock {\em International Journal of Information Technology \& Decision
        Making} {\bf 2018}, {\em 17},~391--466.
    \newblock {\url{https://doi.org/10.1142/S021962201830001X}}.
    
    \bibitem[Sun et~al.(2018)Sun, Dong, and Liu]{Sun2018}
    Sun, L.; Dong, H.; Liu, A.
    \newblock Aggregation functions considering criteria interrelationships in
    fuzzy multi-criteria decision making: State-of-the-art.
    \newblock {\em IEEE Access} {\bf 2018}, {\em 6},~68104--68136.
    \newblock {\url{https://doi.org/10.1109/ACCESS.2018.2879741}}.
    
    \bibitem[Kahraman et~al.(2015)Kahraman, Onar, and Oztaysi]{Kahraman2015}
    Kahraman, C.; Onar, S.C.; Oztaysi, B.
    \newblock Fuzzy multicriteria decision-making: A literature review.
    \newblock {\em International Journal of Computational Intelligence Systems}
    {\bf 2015}, {\em 8},~637--666.
    \newblock {\url{https://doi.org/10.1080/18756891.2015.1046325}}.
    
    \bibitem[Mardani et~al.(2015)Mardani, Jusoh, and Zavadskas]{Mardani2015}
    Mardani, A.; Jusoh, A.; Zavadskas, E.K.
    \newblock Fuzzy multiple criteria decision-making techniques and applications
    -- Two decades review from 1994 to 2014.
    \newblock {\em Expert Systems with Applications} {\bf 2015}, {\em
        42},~4126--4148.
    \newblock {\url{https://doi.org/10.1016/j.eswa.2015.01.003}}.
    
    \bibitem[Petrovic et~al.(2023)Petrovic, Mihajlović, Markovic,
    Hashemkhani~Zolfani, and Madić]{Petrovic2023}
    Petrovic, G.S.; Mihajlović, J.; Markovic, D.; Hashemkhani~Zolfani, S.; Madić,
    M.
    \newblock Comparison of aggregation operators in the group decision-making
    process: A real case study of location selection problem.
    \newblock {\em Sustainability} {\bf 2023}, {\em 15},~8229.
    \newblock {\url{https://doi.org/10.3390/su15108229}}.

    \bibitem[Afzal et~al.(2018)Afzal, Ali, Ali, Hussain, Ali, Khan, Amin, Kang,
    and Lee]{Afzal2018}
    Afzal, M.; Ali, S.I.; Ali, R.; Hussain, M.; Ali, T.; Khan, W.A.; Amin, M.B.;
    Kang, B.H.; Lee, S.
    \newblock Personalization of wellness recommendations using contextual
    interpretation.
    \newblock {\em Expert Systems with Applications} {\bf 2018}, {\em
        96},~506--521.
    \newblock {\url{https://doi.org/10.1016/j.eswa.2017.11.006}}.

    \bibitem[Shyur(2025)]{Shyur2025}
    Shyur, H.J.
    \newblock Combination weighting method using Z-numbers for multi-criteria
    decision-making.
    \newblock {\em Applied Soft Computing} {\bf 2025}, {\em 174},~112992.
    \newblock {\url{https://doi.org/10.1016/j.asoc.2025.112992}}.
    
    \bibitem[P{\'{e}}rez-Ca{\~{n}}edo et~al.(2023)P{\'{e}}rez-Ca{\~{n}}edo, Porras,
    Pelta, and Verdegay]{PerezCanedo2022}
    P{\'{e}}rez-Ca{\~{n}}edo, B.; Porras, C.; Pelta, D.A.; Verdegay, J.L.
    \newblock Modeling Contexts as Fuzzy Propositions in Optimization Problems.
    \newblock {\em {IEEE} Transactions on Fuzzy Systems} {\bf 2023}, {\em
        31},~1474--1483.
    \newblock {\url{https://doi.org/10.1109/tfuzz.2022.3203786}}.
    
    \bibitem[Pérez-Cañedo et~al.(2024)Pérez-Cañedo, Novoa-Hernández, Porras,
    Pelta, and Verdegay]{PerezCanedo2024}
    Pérez-Cañedo, B.; Novoa-Hernández, P.; Porras, C.; Pelta, D.A.; Verdegay,
    J.L.
    \newblock Contextual analysis of solutions in a tourist trip design problem: A
    fuzzy logic-based approach.
    \newblock {\em Applied Soft Computing} {\bf 2024}, {\em 154},~111351.
    \newblock {\url{https://doi.org/10.1016/j.asoc.2024.111351}}.
    
    \bibitem[Zimmermann and Zysno(1980)]{Zimmermann1980}
    Zimmermann, H.; Zysno, P.
    \newblock Latent Connectives in Human Decision Making.
    \newblock {\em Fuzzy Sets and Systems} {\bf 1980}, {\em 4},~37--51.
    \newblock {\url{https://doi.org/10.1016/0165-0114(80)90062-7}}.
    
    \bibitem[Natoli et~al.(2024)Natoli, Feeny, Li, and Zuhair]{Natoli2024}
    Natoli, R.; Feeny, S.; Li, J.; Zuhair, S.
    \newblock Aggregating the Human Development Index: A Non-Compensatory Approach.
    \newblock {\em Social Indicators Research} {\bf 2024}, {\em 172},~499--515.
    \newblock {\url{https://doi.org/10.1007/s11205-024-03318-7}}.
    
    \bibitem[Blancas and Contreras(2024)]{Blancas2024}
    Blancas, F.; Contreras, I.
    \newblock Global {SDG} Composite Indicator: A New Methodological Proposal That
    Combines Compensatory and Non-Compensatory Aggregations.
    \newblock {\em Sustainable Development} {\bf 2024}.
    \newblock {\url{https://doi.org/10.1002/sd.3109}}.

    \bibitem[Teng and Tan(2008)]{TengTan2008}
    Teng, T.H.; Tan, A.H.
    \newblock Cognitive Agents Integrating Rules and Reinforcement Learning for
    Context-Aware Decision Support.
    \newblock In Proceedings of the 2008 IEEE/WIC/ACM International Conference on
    Web Intelligence and Intelligent Agent Technology, 2008, pp. 318--321.
    \newblock {\url{https://doi.org/10.1109/WIIAT.2008.163}}.

    \bibitem[Yao and Kumar(2012)]{YaoKumar2012}
    Yao, W.; Kumar, A.
    \newblock Integrating Clinical Pathways into {CDSS} Using Context and Rules: A
    Case Study in Heart Disease.
    \newblock In Proceedings of the 2nd ACM SIGHIT International Health Informatics
    Symposium, 2012, pp. 267--276.
    \newblock {\url{https://doi.org/10.1145/2110363.2110431}}.

    \bibitem[Cronholm and G{\"o}bel(2022)]{CronholmGoebel2022}
    Cronholm, S.; G{\"o}bel, H.
    \newblock Design Principles for Human-Centred {AI}.
    \newblock In {\em ECIS 2022 Research Papers}; 2022; paper~32.
    \newblock {\url{https://aisel.aisnet.org/ecis2022_rp/32}}.

    \bibitem[Yager(2005)]{Yager2005}
    Yager, R.R.
    \newblock Extending Multicriteria Decision Making by Mixing t-Norms and {OWA}
    Operators.
    \newblock {\em International Journal of Intelligent Systems} {\bf 2005}, {\em
        20},~453--474.
    \newblock {\url{https://doi.org/10.1002/int.20075}}.

    \bibitem[del~Moral et~al.(2017)del~Moral, Chiclana, Tapia~Garc{\'i}a, and
    Herrera-Viedma]{MoralEtAl2017}
    del~Moral, M.J.; Chiclana, F.; Tapia~Garc{\'i}a, J.M.; Herrera-Viedma, E.
    \newblock An Analysis on Consensus Measures in Group Decision Making.
    \newblock In Proceedings of the 2017 4th International Conference on Control,
    Decision and Information Technologies ({CoDIT}), 2017, pp. 0262--0267.
    \newblock {\url{https://doi.org/10.1109/CoDIT.2017.8102605}}.

    \bibitem[Merkert et~al.(2015)Merkert, Mueller, and Hubl]{MerkertEtAl2015}
    Merkert, J.; Mueller, M.; Hubl, M.
    \newblock A Survey of the Application of Machine Learning in Decision Support
    Systems.
    \newblock In Proceedings of the 23rd European Conference on Information Systems
    ({ECIS}), 2015.
    \newblock {\url{https://doi.org/10.18151/7217429}}.

    \bibitem[Csisz{\'a}r(2021)]{Csiszar2021}
    Csisz{\'a}r, O.
    \newblock Ordered Weighted Averaging Operators: A Short Review.
    \newblock {\em IEEE Systems, Man, and Cybernetics Magazine} {\bf 2021}, {\em
        7},~4--12.
    \newblock {\url{https://doi.org/10.1109/MSMC.2020.3036378}}.

    \bibitem[Trillo et~al.(2023)Trillo, Cabrerizo, del~Moral, Morente-Molinera,
    Tapia, and Herrera-Viedma]{TrilloEtAl2023}
    Trillo, J.R.; Cabrerizo, F.J.; del~Moral, M.J.; Morente-Molinera, J.A.;
    Tapia, J.M.; Herrera-Viedma, E.
    \newblock A Consensus-Based Multi-Criteria Group Decision-Making Method Based
    on an Aggregated Operator Customised by Experts.
    \newblock In Proceedings of the 2023 9th International Conference on Control,
    Decision and Information Technologies ({CoDIT}), 2023.
    \newblock {\url{https://doi.org/10.1109/CoDIT58514.2023.10284209}}.

\end{thebibliography}
\PublishersNote{}
\end{adjustwidth}

\end{document}
